\section{Expanded Convergence and Generalization Bridge}
\label{app:expanded-theory-bridge}
All definitions, hypotheses, and proofs are in Appendix~\ref{app:theory-full}.
The central question is how accurately learned gradient subspaces affect an
optimizer whose moments and normalizers also evolve. We follow the chain from
covariance estimation to flag tracking, adaptive-state perturbations, and
parameter descent. For generalization, we examine both sensitivity to one
training example and the complexity of a distribution around the learned
parameters. These routes expose different assumptions and complementary terms.

Write $P_k=U_{:r_k}U_{:r_k}^\top$, $\Pi_k=P_k-P_{k-1}$,
$\Phi(U)=\sum_k\omega_k\operatorname{tr}(P_kC_k)$, and
$d_P^2(U,V)=\sum_k\|P_k(U)-P_k(V)\|_F^2$.
The key geometric reduction is the overlap coupling
$T_{ja}=\operatorname{tr}(\Pi_j(U)\Pi_a(V))$:
Lemma~\ref{lem:rigid-metrics} gives
$d_P^2(U,V)=\sum_{j,a}|j-a|T_{ja}$ exactly. A misplaced direction is charged
only for the prefix boundaries it crosses. This replaces repeated triangle
inequalities and yields the sharp factor $K$ in the assignment-distance
comparison. Hoffman--Kruskal integrality then resolves the commuting optimum.
\begin{gather}
 [C_k,C_l]=0,\quad C_k=Q\operatorname{diag}(\lambda_{ki})Q^\top,\quad
 b_{ji}=\sum_{k\ge j}\omega_k\lambda_{ki},\\
 \max_{[U]}\Phi(U)=
 \max_{\substack{x\ge0,\ \sum_jx_{ji}=1\\\sum_ix_{ji}=q_j}}
           \sum_{j,i}b_{ji}x_{ji}
 =\max_{\substack{x\in\{0,1\}^{(K+1)\times d}\\\sum_jx_{ji}=1,\ \sum_ix_{ji}=q_j}}
           \sum_{j,i}b_{ji}x_{ji}.
 \label{eq:main-assignment}
\end{gather}
Theorem~\ref{thm:tf-assignment} gives~\eqref{eq:main-assignment};
Proposition~\ref{prop:tf-cycle} gives
\begin{equation}
 \exists\{C_k\}:\quad [C_k,C_l]=0,\quad
 \operatorname{grad}\Phi(U_0)=0,\quad\operatorname{Hess}\Phi(U_0)\prec0,
 \quad\Phi(U_0)<\max_U\Phi(U).
\end{equation}
The assignment formula characterizes the global optimum, while the
counterexample shows why local ascent requires a basin assumption even when
the covariance matrices commute. The following tracking bounds therefore
describe a regular local oracle branch.
Under the regular-chart hypotheses of Assumption~\ref{ass:tf-chart},
Lemma~\ref{lem:tf-oracle} and Theorem~\ref{thm:tf-tracking} imply
\begin{align}
 -L_\phi I\preceq\nabla^2f_C\preceq-\mu_\phi I,\quad
 \|\widehat C_k-C_k\|_{\mathrm{op}}\le\epsilon_k
 &\Longrightarrow
 \|x_{\widehat C}^\star-x_C^\star\|
 \le\frac{2J_\psi}{\mu_\phi}\sum_k\omega_k\sqrt{r_k}\epsilon_k,\\
 \sqrt{E_{s+1}}\le q\sqrt{E_s}+b_s,\quad q<1
 &\Longrightarrow
 \sqrt{E_s}\le q^s\sqrt{E_0}+\sum_{j<s}q^{s-1-j}b_j.
\end{align}
Here $b_s$ is the $L^2$ norm of the combined forcing. The optimal quadratic
contraction is $q_\star=(L_\phi-\mu_\phi)/(L_\phi+\mu_\phi)$; centered
noise has variance gain $\xi^2/(1-q^2)$, while persistent drift has RMS gain
$1/(1-q)$.
We first isolate a predictable positive preconditioner and retain the exact
algorithmic discrepancy as an error term. This yields a transparent descent
bound; the analysis of current-gradient normalization below removes that
predictability requirement under explicit noise and state-variation controls.
For $L$-smooth $F$, let $h_t=\nabla F(w_t)$,
$g_t=h_t+b_t+\zeta_t$, $\mathbb E_t\zeta_t=0$,
$\mathbb E_t\|\zeta_t\|^2\le\sigma_t^2$, and
\begin{gather}
 w_{t+1}=w_t-\eta_t(B_tg_t+e_t),\quad
 B_t\in\mathcal F_{t-1},\quad mI\preceq B_t\preceq MI,\\
 \varepsilon_t=e_t+B_tb_t,\quad S_T=\sum_{t<T}\eta_t,\quad
 \mathsf v_t=\operatorname{tr}(B_t^2\Sigma_t).
\end{gather}
Theorems~\ref{thm:tf-convergence} and~\ref{thm:tf-PL} yield
\begin{gather}
 \eta_t\le1/(4LM),\quad\Delta_0:=\mathbb E F(w_0)-F_{\inf},\\
 \tfrac12\|\nabla F\|^2\ge\mu(F-F^\star)\quad\text{(PL case)}.
\end{gather}
\begin{align}
 \frac{\sum_{t<T}\eta_t\mathbb E\|h_t\|^2}{S_T}
 &\le\frac{2\Delta_0}{mS_T}
 +\frac{2L\sum_{t<T}\eta_t^2\mathbb E\mathsf v_t}{mS_T}\\
 &\quad+\frac{2\sum_{t<T}\left[\eta_t\mathbb E\|\varepsilon_t\|_{B_t^{-1}}^2
                 +L\eta_t^2\mathbb E\|\varepsilon_t\|^2\right]}{mS_T},\\
 \mathbb E(F(w_T)-F^\star)
 &\le(1-\mu m\eta)^T\Delta_0
 +\frac{L\eta\overline v+(m^{-1}+L\eta)E^2}{\mu m}.
\end{align}
Here $\Sigma_t=\mathbb E_t[\zeta_t\zeta_t^\top]$,
$\sup_t\mathbb E\mathsf v_t\le\overline v\le M^2\sigma^2$,
$\mathbb E\|\varepsilon_t\|^2\le E^2$, and the PL step-size restrictions
apply. The exact implemented discrepancy is expanded in
Definition~\ref{def:rigid-error} and Proposition~\ref{prop:rigid-transport}.

For neighboring samples $S\sim S^{(i)}$, let
$z_t=(\delta w_t,\delta U_t,\delta Y_t,\delta M_t,\delta\sqrt{V_t})^\top$.
Definition~\ref{def:rigid-transition} and Theorems~\ref{thm:tf-stability}
and~\ref{thm:rigid-generalization} give
\begin{gather}
 z_{t+1}\preceq\mathcal M_tz_t+d_t I_t,\quad z_0=0,\quad
 \mathbb P(I_t=1)=b_t/n,\\
 \left|\mathbb E(F_\mathcal D(w_T)-F_S(w_T))\right|
 \le\epsilon_T^{\mathrm{stab}}
 \le\frac{L_0}{n}e_w\sum_{t<T}b_t
                  (\mathcal M_{T-1}\cdots\mathcal M_{t+1})d_t.
\end{gather}

For $0\le\ell\le1$, a sample-independent prior $P$, and any posterior $Q\ll P$,
Theorem~\ref{thm:rigid-pac} also gives, simultaneously in $Q$ with probability
at least $1-\delta$,
\begin{equation}
 L(Q)\le\operatorname{kl}_+^{-1}\!\left(\widehat L_S(Q),
 \frac{\operatorname{KL}(Q\|P)+\log(\Xi_n/\delta)}n\right).
\end{equation}
Here $\Xi_n$ is the exact Seeger--Maurer moment constant, at most
$2\sqrt n$ for $n\ge8$; retaining inverse KL avoids the Pinsker relaxation.
Proposition~\ref{prop:rigid-pac-flag} evaluates this complexity exactly for a
full-rank Gaussian posterior aligned with the learned flag.

For the implemented current-gradient normalizer $\mathcal T_t(g)=B_t(g)g$,
Theorem~\ref{thm:actual-coercivity} replaces predictability of $B_t$ by an explicit
conditional symmetry defect:
\begin{gather}
 b_t=\mathbb E_tg_t-\nabla F(w_t),\quad
 \delta_t=W_1(\mathcal L(\xi_t\mid\mathcal F_{t-1}),
                    \mathcal L(-\xi_t\mid\mathcal F_{t-1})),\\
 \widehat D_t=\|b_t\|+\|\mathbb E_t\mathcal T_t(\xi_t)\|/M
 \le\|b_t\|+\delta_t/2,\\
 \langle\nabla F(w_t),\mathbb E_t\mathcal T_t(g_t)\rangle
 \ge\tfrac\kappa2\|\nabla F(w_t)\|^2-\tfrac{M^2}{2\kappa}\widehat D_t^2.
\end{gather}
Theorem~\ref{thm:actual-convergence} then controls the actual moment/reset/transport
map by its observable variation budgets. Proposition~\ref{prop:actual-skew-counterexample}
gives a closed-form nonzero stationary bias for asymmetric, unbiased noise.

Under the bounded-loss anisotropic-curvature hypothesis of
Assumption~\ref{ass:directional-curvature}, Theorem~\ref{thm:directional-pac} gives
\begin{gather}
 \mathcal C_\lambda(U)=\frac1{2\lambda}\sum_jmq_j
      \log(1+2\lambda\sigma_0^2\overline h_j(U)),\qquad
 (\sigma_j^2)^\star=\frac{\sigma_0^2}{1+2\lambda\sigma_0^2\overline h_j(U)}.
\end{gather}
Proposition~\ref{prop:directional-refinement} proves a strict decrease in this
optimized posterior cost when a refined block separates unequal average curvatures.
It also bounds the improvement above and below by the weighted variance of
those curvatures. The tracking sensitivity depends on the spectral spread of
the partial trace of $H_0$, so isotropic curvature cancels exactly
(Proposition~\ref{prop:directional-tracking}).
This last comparison identifies when additional posterior directions improve
the bound. Translating it into an optimizer advantage also requires the learned
flag and the resulting parameter trajectory to have the appropriate geometry.
